\chapter{Introduction}
\label{ch:introduction}

Forge is a small, statically typed functional programming language. It is
in the ML family: programs are made of definitions, every expression has a
type, and the compiler works the types out for you. On top of that core it
offers the features that make a functional language pleasant to write ---
algebraic data types and pattern matching, records with structural typing,
Haskell-style \idx{traits} (type classes), Rust-style
\fc|macro_rules!| macros, and modules --- and it runs your programs two ways:
with an \idx{interpreter}, or compiled ahead of time to a native executable
through LLVM.

\begin{itemize}
  \item \textbf{Types without the ceremony.} Hindley--Milner type inference
    with let-polymorphism: annotations are optional and checked when given.
  \item \textbf{Data you can take apart.} Sum types, tuples, records and
    lists, with nested patterns, list comprehensions and warnings for
    non-exhaustive or redundant \fc|case| arms.
  \item \textbf{Overloading done right.} Traits with superclasses, default
    methods, parameterised instances (\fc|impl Show for List<a> requires Show<a>|)
    and higher-kinded classes such as \fc|Functor| and \fc|Monad|.
  \item \textbf{One language, two back ends.} A tree-walking interpreter and
    REPL for quick feedback, and a compiler that turns the same programs into
    native code. The test suite runs thousands of programs through both and
    requires identical output.
\end{itemize}

\section{A guided tour}
\label{sec:tour}

A Forge program is a sequence of top-level definitions. A definition
\fc|let () = |\textit{e} evaluates \textit{e} for its effect; that is how a
program prints.

\forgeexample{tour-hello}

Functions are values. \fc|let f x y = ...| defines a curried function of two
arguments, \fc|fn x -> ...| is an anonymous function, and applying a function
to fewer arguments than it takes gives back a function.

\forgeexample{tour-functions}

\fc|println| prints any value whose type has a \fc|Show| instance --- numbers,
strings, lists, options, tuples --- followed by a newline. Lists come with
ranges and comprehensions, and \fc|case| takes values apart by their shape:

\forgeexample{tour-lists}

New types are declared with \fc|type|. A \emph{sum type} lists its
constructors; records need no declaration at all, and a function that reads a
field accepts any record that has it:

\forgeexample{tour-types}

Traits describe operations that several types support. An instance can be
conditional --- lists of describable things are describable --- and a
function that uses a trait method works for every type with an instance:

\forgeexample{tour-traits}

\section{Running Forge}

The same file runs in the interpreter or compiles to a native executable:

\begin{shellpanel}
\begin{lstlisting}
$ dune exec ./bin/interpreter.exe documentation/examples/tour-hello.forge
Hello, Forge!
$ dune exec ./bin/compile_forge.exe documentation/examples/tour-hello.forge hello
$ ./hello
Hello, Forge!
\end{lstlisting}
\end{shellpanel}

Chapter~\ref{ch:tools} describes the tools; Chapters
\ref{ch:lexical}--\ref{ch:modules} describe the language, Chapter~\ref{ch:formal}
defines it formally, Chapter~\ref{ch:library} documents the standard library and
Chapter~\ref{ch:implementation} explains how the implementation works.

\begin{notebox}[About the examples]
Every program in this manual is a file in \texttt{documentation/examples},
shown together with the output it produces. The test suite
(\texttt{test/manual\_tests.ml}) runs each of them with the interpreter and
with the native compiler and checks that both print exactly what the manual
shows.
\end{notebox}
